\documentclass{article}
\usepackage{amsmath,amssymb}
\title{Why the flat relative-energy comparison does not cover literal torus gluing}\author{}\date{}
\begin{document}\maketitle
Section 6 of \cite{ND} and the all shell relative entropy calculation compare weights of the SAME plaquette words on the SAME field map. Their key estimate $|w|\leq\delta E_0$ is strong and useful, but it cannot be applied directly to the original SV5 gluing difference $\Delta S=S_{\rm big}-S_{\rm per}$, even on flat original small tori.

Use the literal common link variables of SV5 with $q,n\geq3$. Choose an orientation $(\mu,\nu)$ and a singly rewired plaquette at $x_\mu=n-1$, $x_\nu\leq n-2$, in small block $b$. Its new second edge belongs to block $b+e_\mu$ at the wrapped coordinate $y_\mu=0$.

On every small torus put the links in a pure gauge:
\[
 U_e=g_x^{-1}g_{x+\widehat e},
\]
where the endpoint is understood periodically INSIDE that small block. Set every $g_x=I$, except at the vertex $(b+e_\mu,y)$, where set $g_y=H\ne I$. Then every original small-torus plaquette is exactly the identity, so
\[
 S_{\rm per}=0.
\]
The specified new big-torus plaquette contains exactly one nontrivial link from that neighboring block. Its other three links lie in block $b$ and are identities. Its holonomy is $H$ or $H^{-1}$, depending on orientation, hence
\[
 \Delta S=S_{\rm big}\geq s(H)>0.
\]
The matrix $H$ can be arbitrarily close to $I$. Thus this failure persists in arbitrarily small raw link neighborhoods; it is not a large-field example. By continuity it also prevents a uniform relative bound in neighborhoods of this flat configuration. In particular
\[
 |\Delta S|\leq C S_{\rm per}
\]
fails for every finite constant $C$ on the unchanged gluing space.

This proves a precise limit of the new all-shell entropy result: gluing changes how the original links enter plaquette words and how independent block gauge frames are matched. It is not merely a relative perturbation of spatial coefficients on a common field map. Integrating or aligning those gauge frames may be part of a constructive comparison, but it has an actual normalization and boundary interaction; the exact full-cube examples in notes 07--08 show why those terms must remain.

The example does not imply a positive lower bound for the full free-energy defect, because the product Wilson law averages all configurations and the interpolation includes compensating entropy. It is not a counterexample to S94. It identifies the missing map/frame comparison needed before the flat relative-energy estimate can enter a proof of S94.
\begin{thebibliography}{9}
\bibitem{ND} \emph{Nonlinear dispersion and unequal spatial weights}, \texttt{lattice\_theory/block\_and\_transfer/conditional\_estimates/nonlinear\_dispersion\_and\_unequal\_spatial\_weights.tex} in this collection.
\end{thebibliography}
\end{document}
